% =====================================================================
\kicker{Guide topic 2}
\section{GPIO \& UART Integration: System Architecture}\label{sec:gpio}
% =====================================================================
\resultstrip{Connect GPIO and UART to the Stage~2 core and enter the design in ChipInventor.}
  {One decoder change; 16 blocks and 8 Inout Pins on the canvas}
  {Figures~\ref{fig:arch}–\ref{fig:closeup}, Tables~\ref{tab:regmap}–\ref{tab:designdec}}

GPIO and UART are bus slaves behind the address decoder, the only core block modified in Stage~3 (\figref{fig:arch}). The decoder selects a slave from the upper address bits, forwards the bus signals and returns the selected read data to the core. Each peripheral decodes only its word offset, address bits 11:2, as the Block Guide requires~\cite{blockguide}. The firmware therefore reaches the pins and the serial line with ordinary \code{lw} and \code{sw} instructions (\tabref{tab:regmap}); \tabref{tab:designdec} records the design decisions.

\begin{figure}[tbp]
\centering
\includegraphics[width=0.8\linewidth]{fig02_architecture.pdf}
\caption[SoC integration architecture]{SoC integration. The address decoder is the only modified core block; GPIO and UART are new bus slaves, and each pin has its own Inout Pin block. The two firmwares differ only in the ROM.}
\label{fig:arch}
\end{figure}

% Table 1 -- Memory map and peripheral registers (Block Guide, Stage 3; Addr_dec.v, DMEM.v)
\begin{table}[H]
\caption[Memory map and peripheral registers]{Memory map and peripheral registers, from the Block Guide Stage 3 register maps~\cite{blockguide} and \texttt{Addr\_dec.v}.}
\label{tab:regmap}
{\tablestyle
\begin{tabularx}{\linewidth}{@{}P{27mm}P{27mm}L@{}}
\toprule
\colhead{Block} & \colhead{Base address} & \colhead{Registers (offset) / contents}\\
\midrule
IMEM (ROM) & \hex{0x0040\_0000} & \code{IMEM\_Data\_Center} (275 words) or the official test firmware (21 words)\\
DMEM & \hex{0x1001\_0000} & 8\,KB SRAM (2,048 words); application state in 22 words, cleared at boot\\
GPIO & \hex{0xF000\_0000} & \code{DATAOUT} (\hex{0x0}), \code{DATAIN} (\hex{0x4}), \code{DATADIR} (\hex{0x8}; 1 = output)\\
GPIO\_PIN\_SPLIT & none (wiring) & none; per pin \code{C} = \code{DATAOUT}[$k$], \code{D} = \code{DATADIR}[$k$], \code{P} = pin\\
UART & \hex{0xF100\_0000} & \code{TXDATA} (\hex{0x0}), \code{RXDATA} (\hex{0x4}), \code{CONTROL} (\hex{0x8}): \code{TRANSMIT} bit~0, \code{RXDONE} bit~1, \code{TXDONE} bit~2\\
\bottomrule
\end{tabularx}}
\end{table}

% Table 2 -- Integration design decisions (GPIO.v, GPIO_PIN_SPLIT.v, UART.v, top.v)
\begin{table}[H]
\caption[Integration design decisions]{Integration design decisions and their rationale.}
\label{tab:designdec}
{\tablestyle
\begin{tabularx}{\linewidth}{@{}P{22mm}P{68mm}L@{}}
\toprule
\colhead{Area} & \colhead{Decision} & \colhead{Rationale}\\
\midrule
GPIO reset & All pins are inputs; \code{DATAOUT} = 0 & No pin is driven before the firmware writes \code{DATADIR}.\\
GPIO input & 2-flop synchronizer before \code{DATAIN}; in output mode a pin's \code{DATAIN} bit holds its last input value & Pins are asynchronous to the clock; the Guide isolates the input buffer of an output pin (§1.3.1)~\cite{blockguide}.\\
Pin tristate & \code{GPIO\_PIN\_SPLIT} feeds eight 1-bit Inout Pin blocks: \code{pins\_io\_$k$ = D ? C : 1\sq bz} & An 8-bit Inout Pin applies one enable to all pins, and the export drops a custom block's \code{inout} port; scalar pins avoid both without editing \code{hdl.v}.\\
UART baud & $N = 25\ \mathrm{MHz} / 115\,200 = 217$ clocks per bit & Integer divisor of the system clock; 0.01\,\% baud error.\\
UART transmit & A frame starts only when \code{TRANSMIT} is set (\code{TX\_ON\_TXDATA\_WRITE = 0}); \code{TRANSMIT} self-clears; \code{TXDONE} = 1 when idle & Specified behaviour: writing \code{TXDATA} only loads the byte.\\
UART receive & Synchronized input; start bit re-checked at half a bit; one sample per bit centre; frame dropped if the stop bit is low; \code{RXDONE} cleared by writing 0, and a byte arriving in the same cycle takes priority & Rejects glitches and framing errors without oversampling; no received byte is lost at the clear.\\
\bottomrule
\end{tabularx}}
\end{table}


\subsection{The ChipInventor block diagram}\label{sec:canvas}
\figref{fig:canvas} is the final design as entered on the canvas, with one wire colour per bus; the numbered markers identify the blocks. The canvas exports \code{hdl.v}, submitted as \code{rtl/hdl/hdl\_application.v} and used by every test in \secref{sec:validation} without modification. Each custom block is identical to the repository RTL.

\medskip
\noindent\begin{minipage}[t]{0.715\linewidth}
\vspace{0pt}
{\setlength{\fboxsep}{0pt}\color{hair}\fbox{%
\begin{tikzpicture}
  \node[inner sep=0, anchor=south west] (img) at (0,0)
       {\includegraphics[width=\dimexpr\linewidth-2\fboxrule\relax]{fig03_chipinventor_canvas.jpg}};
  \begin{scope}[x={(img.south east)}, y={(img.north west)}]
    % marker positions: fractions of width (x) and height from the bottom (y)
    \foreach \x/\y/\n in {0.605/0.620/1, 0.820/0.452/2, 0.131/0.670/3, 0.250/0.724/4,
                          0.251/0.539/5, 0.246/0.447/6, 0.341/0.422/7, 0.431/0.488/8, 0.245/0.265/9}
      \node[circle, fill=alarm, text=white, inner sep=0pt, minimum size=4.2mm,
            font=\sffamily\bfseries\scriptsize, draw=white, line width=0.6pt] at (\x,\y) {\n};
  \end{scope}
\end{tikzpicture}}}
\end{minipage}\hfill
\begin{minipage}[t]{0.265\linewidth}
\vspace{0pt}\sidecaption
\captionof{figure}[Block diagram generated in ChipInventor (mandatory)]{Block diagram generated in ChipInventor (mandatory).}
\label{fig:canvas}
\vspace{2pt}
{\sffamily\footnotesize\renewcommand{\arraystretch}{1.12}%
\begin{tabular}{@{}l@{\hspace{4pt}}p{0.8\linewidth}@{}}
\mk{1} & Control Unit\\
\mk{2} & PC, ALU, MULT, CRC16, Register File, LSU\\
\mk{3} & Address decoder\\
\mk{4} & IMEM ROM, 275 words\\
\mk{5} & DMEM, 8\,KB\\
\mk{6} & GPIO\\
\mk{7} & \code{GPIO\_PIN\_SPLIT}\\
\mk{8} & Inout Pins \code{pins\_io\_0}–\code{7}\\
\mk{9} & UART
\end{tabular}}
\end{minipage}

\medskip
\noindent\begin{minipage}[b]{0.56\linewidth}
{\setlength{\fboxsep}{0pt}\color{hair}\fbox{%
\begin{tikzpicture}
  \node[inner sep=0, anchor=south west] (img) at (0,0)
       {\includegraphics[width=\dimexpr\linewidth-2\fboxrule\relax]{fig04_canvas_closeup.jpg}};
  \begin{scope}[x={(img.south east)}, y={(img.north west)}]
    \foreach \x/\y/\n in {0.292/0.858/6, 0.553/0.779/7, 0.800/0.950/8, 0.287/0.297/9}
      \node[circle, fill=alarm, text=white, inner sep=0pt, minimum size=4.2mm,
            font=\sffamily\bfseries\scriptsize, draw=white, line width=0.6pt] at (\x,\y) {\n};
  \end{scope}
\end{tikzpicture}}}
\end{minipage}\hfill
\begin{minipage}[b]{0.415\linewidth}\sidecaption
\captionof{figure}[Close-up of the peripheral side of the canvas]{Close-up of \figref{fig:canvas}. GPIO \mk{6} sends \texttt{o\_Pins\_Out} and \texttt{o\_Pins\_Oe} to \texttt{GPIO\_PIN\_SPLIT} \mk{7}, which drives \texttt{C} and \texttt{D} of each Inout Pin \mk{8} and returns each pin on \texttt{i\_P0}–\texttt{i\_P7}. The UART \mk{9} connects to \texttt{rx\_i} and \texttt{tx\_o} with \texttt{TX\_ON\_TXDATA\_WRITE = 0}.}
\label{fig:closeup}
\end{minipage}
